%% Generated by tools/gen_error_codes.py from docs/errors/ of the compiler;
%% edit the compiler's pages or tools/error_themes.txt, then rerun it.

\clearpage
\section{Syntax}
\label{sec:codes:syntax}

\begin{longtable}{@{}l p{0.78\linewidth} r@{}}
\toprule Code & Message & Page\\ \midrule \endhead
\hyperref[sec:codes:e2001]{E2001} & block is never closed & \pageref{sec:codes:e2001}\\
\hyperref[sec:codes:e2002]{E2002} & decorator \errhole{} is defined multiple times & \pageref{sec:codes:e2002}\\
\hyperref[sec:codes:e2004]{E2004} & scope guard \errhole{} is declared multiple times & \pageref{sec:codes:e2004}\\
\hyperref[sec:codes:e2006]{E2006} & decorator \errhole{} cannot be applied to a subpattern variable & \pageref{sec:codes:e2006}\\
\hyperref[sec:codes:e2007]{E2007} & unexpected \errhole{} & \pageref{sec:codes:e2007}\\
\hyperref[sec:codes:e2008]{E2008} & unexpected attribute list \errhole{} & \pageref{sec:codes:e2008}\\
\hyperref[sec:codes:e2009]{E2009} & read \errhole{}, but expected \errhole{} & \pageref{sec:codes:e2009}\\
\hyperref[sec:codes:e2012]{E2012} & a named move ctor requires an argument list & \pageref{sec:codes:e2012}\\
\hyperref[sec:codes:e2013]{E2013} & missing the value moved by the operator & \pageref{sec:codes:e2013}\\
\hyperref[sec:codes:e2014]{E2014} & attribute \errhole{} takes no argument & \pageref{sec:codes:e2014}\\
\hyperref[sec:codes:e2015]{E2015} & attribute \errhole{} takes \errhole{} argument(s), but \errhole{} were given & \pageref{sec:codes:e2015}\\
\bottomrule
\end{longtable}

\errorcode{E2001}{block is never closed}
\label{sec:codes:e2001}

Raised during \textbf{parsing}, from \texttt{SyntaxErrorMessage::\allowbreak{}BLOCK\_\allowbreak{}NEVER\_\allowbreak{}CLOSED} (\textit{src/\allowbreak{}ymirc/\allowbreak{}syntax/\allowbreak{}errors.yr}).

\subsubsection*{What it means}

The parser reached the end of the file, or a token that cannot appear there, while a block opened earlier was still waiting for its closing brace. The diagnostic underlines the whole of what the block swallowed, from the opening brace to the end of the file: the imbalance starts at the brace, and everything after it was read as part of the block.

\subsubsection*{Example}

\textit{test\_\allowbreak{}resources/\allowbreak{}diagnostics/\allowbreak{}test87.yr}:

%% check: skip
\begin{lstlisting}[style=coloredverbatimError]
// a function body never closed: the region runs from the '{' to the end of the file
fn main () {
    let a = 12;
    foo (a);
\end{lstlisting}

\begin{lstlisting}[style=bashVerb, breaklines=true, escapechar=~]
Error[E2001] : block is never closed
 --> test_resources/diagnostics/test87.yr:(2,12)
 2  ~\lstbox{\textSFxi{}}~   fn main () {
    ~\lstbox{\textSFv{}}~ ~\lstbox{\textSFi{}}~~\lstbox{\textSFx{}}~~\lstbox{\textSFx{}}~~\lstbox{\textSFx{}}~~\lstbox{\textSFx{}}~~\lstbox{\textSFx{}}~~\lstbox{\textSFx{}}~~\lstbox{\textSFx{}}~~\lstbox{\textSFx{}}~~\lstbox{\textSFx{}}~~\lstbox{\textSFx{}}~~\lstbox{\textSFx{}}~~\lstbox{\textSFx{}}~^
 3  ~\lstbox{\textSFxi{}}~ ~\lstbox{\textSFxi{}}~     let a = 12;
 4  ~\lstbox{\textSFxi{}}~ ~\lstbox{\textSFxi{}}~     foo (a);
    ~\lstbox{\textSFv{}}~ ~\lstbox{\textSFii{}}~~\lstbox{\textSFx{}}~~\lstbox{\textSFx{}}~~\lstbox{\textSFx{}}~~\lstbox{\textSFx{}}~~\lstbox{\textSFx{}}~~\lstbox{\textSFx{}}~~\lstbox{\textSFx{}}~~\lstbox{\textSFx{}}~~\lstbox{\textSFx{}}~~\lstbox{\textSFx{}}~~\lstbox{\textSFx{}}~~\lstbox{\textSFx{}}~^
    ~\lstbox{\textSFii{}}~~\lstbox{\textSFx{}}~~\lstbox{\textSFx{}}~~\lstbox{\textSFx{}}~~\lstbox{\textSFx{}}~~\lstbox{\textSFx{}}~~\lstbox{\textSFx{}}~~\lstbox{\textSFx{}}~
Error[E2009] : read (, but expected {
 --> test_resources/diagnostics/test87.yr:(2,9)
 2  ~\lstbox{\textSFxi{}}~ fn main () {
    ~\lstbox{\textSFv{}}~         ^  note: when reading template symbol
    ~\lstbox{\textSFii{}}~~\lstbox{\textSFx{}}~~\lstbox{\textSFx{}}~~\lstbox{\textSFx{}}~~\lstbox{\textSFx{}}~~\lstbox{\textSFx{}}~~\lstbox{\textSFx{}}~~\lstbox{\textSFx{}}~
\end{lstlisting}

\subsubsection*{How to fix it}

Close the block. When the braces look balanced, the culprit is usually a brace opened inside a string literal or a comment further up, or a missing one on an earlier declaration that swallowed the rest of the file.

\errorcode{E2002}{decorator \errhole{} is defined multiple times}
\label{sec:codes:e2002}

Raised during \textbf{parsing}, from \texttt{SyntaxErrorMessage::\allowbreak{}MULTIPLE\_\allowbreak{}ATTRS} (\textit{src/\allowbreak{}ymirc/\allowbreak{}syntax/\allowbreak{}errors.yr}).

\subsubsection*{What it means}

The same decorator was written twice on one declaration. A decorator states a property; stating it twice says nothing more and is more likely a mistake than an intention.

\subsubsection*{Example}

\textit{test\_\allowbreak{}resources/\allowbreak{}diagnostics/\allowbreak{}test1.yr}:

%% check: skip
\begin{lstlisting}[style=coloredverbatimError]
@final @final
fn foo ()-> i32 { 0 }

fn main () { foo (); }
\end{lstlisting}

\begin{lstlisting}[style=bashVerb, breaklines=true, escapechar=~]
Error[E2002] : decorator final is defined multiple times
 --> test_resources/diagnostics/test1.yr:(1,2)
 1  ~\lstbox{\textSFxi{}}~ @final @final
    ~\lstbox{\textSFv{}}~  ^^^^^  ^^^^^
    ~\lstbox{\textSFii{}}~~\lstbox{\textSFx{}}~~\lstbox{\textSFx{}}~~\lstbox{\textSFx{}}~~\lstbox{\textSFx{}}~~\lstbox{\textSFx{}}~~\lstbox{\textSFx{}}~~\lstbox{\textSFx{}}~
\end{lstlisting}

\subsubsection*{How to fix it}

Remove the duplicate. If two different decorators were meant, check the spelling of the second one.

\errorcode{E2004}{scope guard \errhole{} is declared multiple times}
\label{sec:codes:e2004}

Raised during \textbf{parsing}, from \texttt{SyntaxErrorMessage::\allowbreak{}MULTIPLE\_\allowbreak{}SAME\_\allowbreak{}GUARD} (\textit{src/\allowbreak{}ymirc/\allowbreak{}syntax/\allowbreak{}errors.yr}).

\subsubsection*{What it means}

The same kind of scope guard was declared twice in one block. A block has at most one \tokennolst{exit}, one \tokennolst{success} and one \tokennolst{failure} guard.

\subsubsection*{Example}

\textit{test\_\allowbreak{}resources/\allowbreak{}diagnostics/\allowbreak{}test2.yr}:

%% check: skip
\begin{lstlisting}[style=coloredverbatimError]
fn main () {
    {
    } exit {
    } exit {
    }
}
\end{lstlisting}

\begin{lstlisting}[style=bashVerb, breaklines=true, escapechar=~]
Error[E2004] : scope guard exit is declared multiple times
 --> test_resources/diagnostics/test2.yr:(3,7)
 3  ~\lstbox{\textSFxi{}}~     } exit {
    ~\lstbox{\textSFv{}}~       ^^^^
   ...
 --> test_resources/diagnostics/test2.yr:(4,7)
 4  ~\lstbox{\textSFxi{}}~     } exit {
    ~\lstbox{\textSFv{}}~       ^^^^
    ~\lstbox{\textSFii{}}~~\lstbox{\textSFx{}}~~\lstbox{\textSFx{}}~~\lstbox{\textSFx{}}~~\lstbox{\textSFx{}}~~\lstbox{\textSFx{}}~~\lstbox{\textSFx{}}~~\lstbox{\textSFx{}}~
\end{lstlisting}

\subsubsection*{How to fix it}

Merge the two bodies into a single guard. Two guards of the same kind would have no defined order between them, which is why the second is rejected rather than chained.

\errorcode{E2006}{decorator \errhole{} cannot be applied to a subpattern variable}
\label{sec:codes:e2006}

Raised during \textbf{parsing}, from \texttt{SyntaxErrorMessage::\allowbreak{}UNDEF\_\allowbreak{}DECORATOR\_\allowbreak{}IN\_\allowbreak{}MATCH} (\textit{src/\allowbreak{}ymirc/\allowbreak{}syntax/\allowbreak{}errors.yr}).

\subsubsection*{What it means}

A decorator was written on a variable of a subpattern inside a pattern match. The subpattern variables of a match are bound by the pattern itself; they take their mutability and their storage from the value being matched, so a decorator on them has nothing to apply to.

\subsubsection*{Example}

\textit{test\_\allowbreak{}resources/\allowbreak{}lazyness/\allowbreak{}test14.yr}:

%% check: skip
\begin{lstlisting}[style=coloredverbatimError]
class A {
    pub let x : i32 = 0;
    pub self () {}
}

fn main () {
    let lazy  A(x-> x) = lazy A();
}
\end{lstlisting}

\begin{lstlisting}[style=bashVerb, breaklines=true, escapechar=~]
Error[E2006] : decorator lazy cannot be applied to a subpattern variable
 --> test_resources/lazyness/test14.yr:(7,9)
 7  ~\lstbox{\textSFxi{}}~     let lazy  A(x-> x) = lazy A();
    ~\lstbox{\textSFv{}}~         ^^^^   ^  note: failed to read pattern variable declaration
    ~\lstbox{\textSFii{}}~~\lstbox{\textSFx{}}~~\lstbox{\textSFx{}}~~\lstbox{\textSFx{}}~~\lstbox{\textSFx{}}~~\lstbox{\textSFx{}}~~\lstbox{\textSFx{}}~~\lstbox{\textSFx{}}~
\end{lstlisting}

\subsubsection*{How to fix it}

Remove the decorator. When the binding needs to be mutable, the decorator belongs on the outermost variable of the pattern, not on the nested one.

\errorcode{E2007}{unexpected \errhole{}}
\label{sec:codes:e2007}

Raised during \textbf{parsing}, from \texttt{SyntaxErrorMessage::\allowbreak{}UNEXPECTED} (\textit{src/\allowbreak{}ymirc/\allowbreak{}syntax/\allowbreak{}errors.yr}).

\subsubsection*{What it means}

The parser read a token that cannot appear at that point of the grammar and had no more specific expectation to report. The token itself is in the message, a file stopping there naming itself ``end of file''.

\subsubsection*{Example}

\textit{test\_\allowbreak{}resources/\allowbreak{}diagnostics/\allowbreak{}test3.yr}:

%% check: skip
\begin{lstlisting}[style=coloredverbatimError]
class Foo {
    @final let x: i32;
    pub self (x: i32) with x = x {}
}

fn main () {}
\end{lstlisting}

\begin{lstlisting}[style=bashVerb, breaklines=true, escapechar=~]
Error[E2007] : unexpected @
 --> test_resources/diagnostics/test3.yr:(2,5)
 2  ~\lstbox{\textSFxi{}}~     @final let x: i32;
    ~\lstbox{\textSFv{}}~     ^  note: when reading template symbol
    ~\lstbox{\textSFii{}}~~\lstbox{\textSFx{}}~~\lstbox{\textSFx{}}~~\lstbox{\textSFx{}}~~\lstbox{\textSFx{}}~~\lstbox{\textSFx{}}~~\lstbox{\textSFx{}}~~\lstbox{\textSFx{}}~
\end{lstlisting}

\subsubsection*{How to fix it}

Read the reported token in its context: this is usually a missing separator, a missing closing delimiter earlier on, or a keyword used where an expression was expected.

\errorcode{E2008}{unexpected attribute list \errhole{}}
\label{sec:codes:e2008}

Raised during \textbf{parsing}, from \texttt{SyntaxErrorMessage::\allowbreak{}UNEXPECTED\_\allowbreak{}ATTRIBUTES} (\textit{src/\allowbreak{}ymirc/\allowbreak{}syntax/\allowbreak{}errors.yr}).

\subsubsection*{What it means}

An attribute list was written on a declaration that does not take one.

\subsubsection*{Example}

\textit{test\_\allowbreak{}resources/\allowbreak{}diagnostics/\allowbreak{}test3.yr}:

%% check: skip
\begin{lstlisting}[style=coloredverbatimError]
class Foo {
    @final let x: i32;
    pub self (x: i32) with x = x {}
}

fn main () {}
\end{lstlisting}

\begin{lstlisting}[style=bashVerb, breaklines=true, escapechar=~]
Error[E2008] : unexpected attribute list {final}
 --> test_resources/diagnostics/test3.yr:(2,5)
 2  ~\lstbox{\textSFxi{}}~     @final let x: i32;
    ~\lstbox{\textSFv{}}~     ^
    ~\lstbox{\textSFxi{}}~ Note :
    ~\lstbox{\textSFxi{}}~  --> test_resources/diagnostics/test3.yr:(2,12)
    ~\lstbox{\textSFxi{}}~  2  ~\lstbox{\textSFxi{}}~     @final let x: i32;
    ~\lstbox{\textSFxi{}}~     ~\lstbox{\textSFv{}}~            ^^^
    ~\lstbox{\textSFxi{}}~     ~\lstbox{\textSFii{}}~~\lstbox{\textSFx{}}~~\lstbox{\textSFx{}}~~\lstbox{\textSFx{}}~~\lstbox{\textSFx{}}~~\lstbox{\textSFx{}}~~\lstbox{\textSFx{}}~~\lstbox{\textSFx{}}~
    ~\lstbox{\textSFii{}}~~\lstbox{\textSFx{}}~~\lstbox{\textSFx{}}~~\lstbox{\textSFx{}}~~\lstbox{\textSFx{}}~~\lstbox{\textSFx{}}~~\lstbox{\textSFvii{}}~~\lstbox{\textSFx{}}~
\end{lstlisting}

\subsubsection*{How to fix it}

Remove the list, or move it onto the declaration it was meant for. Attribute lists apply to declarations, not to the statements inside them.

\errorcode{E2009}{read \errhole{}, but expected \errhole{}}
\label{sec:codes:e2009}

Raised during \textbf{parsing}, from \texttt{SyntaxErrorMessage::\allowbreak{}UNEXPECTED\_\allowbreak{}BUT\_\allowbreak{}LST} (\textit{src/\allowbreak{}ymirc/\allowbreak{}syntax/\allowbreak{}errors.yr}).

\subsubsection*{What it means}

The parser expected one of a known set of tokens and read something else. Both what was read and what was expected are in the message, so the set of possibilities at that point is explicit. A file stopping before any of them is read names itself ``end of file''.

\subsubsection*{Example}

\textit{test\_\allowbreak{}resources/\allowbreak{}lit\_\allowbreak{}class/\allowbreak{}test7.yr}:

%% check: skip
\begin{lstlisting}[style=coloredverbatimError]
class Z {
    pub self ()-> &Z {}
}
\end{lstlisting}

\begin{lstlisting}[style=bashVerb, breaklines=true, escapechar=~]
Error[E2009] : read ->, but expected {
 --> test_resources/lit_class/test7.yr:(2,16)
 2  ~\lstbox{\textSFxi{}}~     pub self ()-> &Z {}
    ~\lstbox{\textSFv{}}~                ^^
    ~\lstbox{\textSFii{}}~~\lstbox{\textSFx{}}~~\lstbox{\textSFx{}}~~\lstbox{\textSFx{}}~~\lstbox{\textSFx{}}~~\lstbox{\textSFx{}}~~\lstbox{\textSFx{}}~~\lstbox{\textSFx{}}~
\end{lstlisting}

\subsubsection*{How to fix it}

Replace the token with one of the listed alternatives. A long list of expectations usually means the construct itself was mis-started a few tokens earlier.

\errorcode{E2012}{a named move ctor requires an argument list}
\label{sec:codes:e2012}

Raised during \textbf{parsing}, from \texttt{SyntaxErrorMessage::\allowbreak{}MOVE\_\allowbreak{}NAMED\_\allowbreak{}CTOR\_\allowbreak{}PARAMS} (\textit{src/\allowbreak{}ymirc/\allowbreak{}syntax/\allowbreak{}errors.yr}).

\subsubsection*{What it means}

A \tokennolst{move} naming the ctor that resets the moved value was not followed by an argument list. The parentheses are what separate the arguments of the ctor from the value being moved, so they are mandatory as soon as a name is given -- even when the ctor takes nothing.

\subsubsection*{Example}

\textit{test\_\allowbreak{}resources/\allowbreak{}diagnostics/\allowbreak{}test90.yr}:

%% check: skip
\begin{lstlisting}[style=coloredverbatimError]
entity Foo {
    pub self open () {}
}

fn main () {
    let mut a = Foo::open ();
    let mut _b_ = move::open a;
}
\end{lstlisting}

\begin{lstlisting}[style=bashVerb, breaklines=true, escapechar=~]
Error[E2012] : a named move ctor requires an argument list
 --> test_resources/diagnostics/test90.yr:(8,23)
 8  ~\lstbox{\textSFxi{}}~     let mut _b_ = move::open a;
    ~\lstbox{\textSFv{}}~                       ^^----  note: failed to read standard variable declaration
    ~\lstbox{\textSFii{}}~~\lstbox{\textSFx{}}~~\lstbox{\textSFx{}}~~\lstbox{\textSFx{}}~~\lstbox{\textSFx{}}~~\lstbox{\textSFx{}}~~\lstbox{\textSFx{}}~~\lstbox{\textSFx{}}~
\end{lstlisting}

\subsubsection*{How to fix it}

Write the empty argument list: \tokennolst{move::open () a}.

\errorcode{E2013}{missing the value moved by the operator}
\label{sec:codes:e2013}

Raised during \textbf{parsing}, from \texttt{SyntaxErrorMessage::\allowbreak{}MOVE\_\allowbreak{}MISSING\_\allowbreak{}OPERAND} (\textit{src/\allowbreak{}ymirc/\allowbreak{}syntax/\allowbreak{}errors.yr}).

\subsubsection*{What it means}

A \tokennolst{(} directly following \tokennolst{move} always opens the argument list of the reset ctor, so what looks like the moved value is read as an argument and the operator is left without an operand.

\subsubsection*{Example}

\textit{test\_\allowbreak{}resources/\allowbreak{}diagnostics/\allowbreak{}test89.yr}:

%% check: skip
\begin{lstlisting}[style=coloredverbatimError]
entity Foo {
    pub self () {}
}

fn foo ()-> Foo;

fn main () {
    let mut _a_ = move (foo ());
}
\end{lstlisting}

\begin{lstlisting}[style=bashVerb, breaklines=true, escapechar=~]
Error[E2013] : missing the value moved by the operator
 --> test_resources/diagnostics/test89.yr:(9,19)
 9  ~\lstbox{\textSFxi{}}~     let mut _a_ = move (foo ());
    ~\lstbox{\textSFv{}}~                   ^^^^ -       -
    ~\lstbox{\textSFxi{}}~                      ~\lstbox{\textSFii{}}~~\lstbox{\textSFx{}}~ note: failed to read standard variable declaration
    ~\lstbox{\textSFxi{}}~                        ~\lstbox{\textSFii{}}~~\lstbox{\textSFx{}}~ note: the argument list of move was read instead, write move () (...) to move a parenthesized value
    ~\lstbox{\textSFxi{}}~                                ~\lstbox{\textSFii{}}~~\lstbox{\textSFx{}}~ error: unexpected ;
    ~\lstbox{\textSFii{}}~~\lstbox{\textSFx{}}~~\lstbox{\textSFx{}}~~\lstbox{\textSFx{}}~~\lstbox{\textSFx{}}~~\lstbox{\textSFx{}}~~\lstbox{\textSFx{}}~~\lstbox{\textSFx{}}~
\end{lstlisting}

\subsubsection*{How to fix it}

Drop the parentheses around the moved value, or -- when they are needed to enclose a complex expression -- put an explicit argument list in front of them: \tokennolst{move () (foo ())}.

\errorcode{E2014}{attribute \errhole{} takes no argument}
\label{sec:codes:e2014}

Raised during \textbf{parsing}, from \texttt{SyntaxErrorMessage::\allowbreak{}ATTRIBUTE\_\allowbreak{}WITHOUT\_\allowbreak{}ARGS} (\textit{src/\allowbreak{}ymirc/\allowbreak{}syntax/\allowbreak{}errors.yr}).

\subsubsection*{What it means}

An attribute declares how many arguments it takes, and most take none. The set of attributes is closed, so this is checked as the name is read, next to the error raised for an attribute that does not exist at all.

Only \tokennolst{@parameters} takes an argument today.

\subsubsection*{Example}

\textit{test\_\allowbreak{}resources/\allowbreak{}diagnostics/\allowbreak{}test91.yr}:

%% check: skip
\begin{lstlisting}[style=coloredverbatimError]
@final(3)
fn main () {}
\end{lstlisting}

\begin{lstlisting}[style=bashVerb, breaklines=true, escapechar=~]
Error[E2014] : attribute final takes no argument
 --> test_resources/diagnostics/test91.yr:(2,2)
 2  ~\lstbox{\textSFxi{}}~ @final(3)
    ~\lstbox{\textSFv{}}~  ^^^^^
    ~\lstbox{\textSFii{}}~~\lstbox{\textSFx{}}~~\lstbox{\textSFx{}}~~\lstbox{\textSFx{}}~~\lstbox{\textSFx{}}~~\lstbox{\textSFx{}}~~\lstbox{\textSFx{}}~~\lstbox{\textSFx{}}~
\end{lstlisting}

\subsubsection*{How to fix it}

Drop the argument list.

\errorcode{E2015}{attribute \errhole{} takes \errhole{} argument(s), but \errhole{} were given}
\label{sec:codes:e2015}

Raised during \textbf{parsing}, from \texttt{SyntaxErrorMessage::\allowbreak{}ATTRIBUTE\_\allowbreak{}ARITY} (\textit{src/\allowbreak{}ymirc/\allowbreak{}syntax/\allowbreak{}errors.yr}).

\subsubsection*{What it means}

An attribute taking arguments takes a fixed number of them, and was written with a different number -- or, as below, with none at all.

\subsubsection*{Example}

\textit{test\_\allowbreak{}resources/\allowbreak{}diagnostics/\allowbreak{}test92.yr}:

%% check: skip
\begin{lstlisting}[style=coloredverbatimError]
@parameters
__test (value: i32) {
    assert (value > 0);
}
\end{lstlisting}

\begin{lstlisting}[style=bashVerb, breaklines=true, escapechar=~]
Error[E2015] : attribute parameters takes 1 argument(s), but 0 were given
 --> test_resources/diagnostics/test92.yr:(2,2)
 2  ~\lstbox{\textSFxi{}}~ @parameters
    ~\lstbox{\textSFv{}}~  ^^^^^^^^^^
    ~\lstbox{\textSFii{}}~~\lstbox{\textSFx{}}~~\lstbox{\textSFx{}}~~\lstbox{\textSFx{}}~~\lstbox{\textSFx{}}~~\lstbox{\textSFx{}}~~\lstbox{\textSFx{}}~~\lstbox{\textSFx{}}~
\end{lstlisting}

\subsubsection*{How to fix it}

Give the attribute the arguments it takes. \tokennolst{@parameters} takes the one value producing the parameter sets of the test.
